\documentclass[lettersize,journal]{IEEEtran}
\usepackage{amsmath,amsfonts}
\usepackage{algorithmic}
\usepackage{algorithm}
\usepackage{array}
\usepackage[caption=false,font=normalsize,labelfont=sf,textfont=sf]{subfig}
\usepackage{textcomp}
\usepackage{stfloats}
\usepackage{url}
\usepackage{verbatim}
\usepackage{graphicx}
\usepackage{cite}
\usepackage{booktabs}
\usepackage{multirow}
\usepackage{xcolor}

\begin{document}

\title{EdgeVLA: Fast and Lightweight Vision-Language-Action Policy with Edge NPU Acceleration for Real-Time Robot Manipulation}

\author{Chunyang Li, and Collaborators% <-this % stops a space
\thanks{Manuscript submitted to IEEE Robotics and Automation Letters (RA-L) / IROS.}
}

\markboth{IEEE Robotics and Automation Letters}{Li \MakeLowercase{\textit{et al.}}: EdgeVLA: Fast and Lightweight VLA Policy on Embedded NPUs}

\maketitle

\begin{abstract}
Vision-Language-Action (VLA) foundation models demonstrate remarkable cross-task generalization in robot manipulation. However, existing state-of-the-art VLA models (e.g., OpenVLA-7B, RT-2-55B) suffer from massive computational footprints ($>7\text{B}$ parameters) and excessive inference latencies ($>500\text{ ms}$), rendering them impractical for real-time closed-loop control on resource-constrained embedded edge robots. In this paper, we propose \textbf{EdgeVLA}, an open-source, fast, and data-efficient Vision-Language-Action policy framework specifically optimized for embedded Neural Processing Units (NPUs). EdgeVLA features: (1) a compact multi-modal architecture with inverted-residual vision extraction and efficient linguistic projection, reducing parameters to $12.4\text{M}$ ($>500\times$ smaller than OpenVLA); (2) an Action Chunking decoder predicting multi-step smooth trajectories ($\mathbf{A} \in \mathbb{R}^{K \times 7}$), effectively reducing query overhead while ensuring physical motion smoothness; and (3) a full-stack hardware-aware compilation pipeline (PyTorch $\to$ ONNX $\to$ Ascend ATC) targeting the Orange Pi AIpro ($20\text{ TOPS}$ Ascend 310B NPU). Extensive evaluations on the standard \textbf{LIBERO} robot manipulation benchmark demonstrate that EdgeVLA achieves $73.4\%$ multi-task success rate while delivering an ultra-low inference latency of $18.2\text{ ms}$ ($54.9\text{ Hz}$ control frequency) on the edge NPU, enabling real-time, low-cost embodied robotic control.
\end{abstract}

\begin{IEEEkeywords}
Vision-Language-Action (VLA), Embodied AI, Embedded NPU, Real-Time Robot Control, Model Compression.
\end{IEEEkeywords}

\section{Introduction}
\IEEEPARstart{V}{ision-Language-Action} (VLA) foundation models have recently emerged as a transformative paradigm for robotic manipulation \cite{brohan2023rt2, kim2024openvla, octo_2023}. By bridging large-scale vision-language pre-training with robotic motor control, VLAs enable robots to follow diverse natural language instructions and generalize to novel tasks.

Despite their promising generalization capabilities, deploying current VLA models onto real-world embedded robotic platforms faces significant bottlenecks:
\begin{enumerate}
    \item \textbf{Excessive Computational Overhead:} State-of-the-art models like OpenVLA-7B \cite{kim2024openvla} require high-end datacenter GPUs (e.g., NVIDIA A100/RTX 4090) and exceed $14\text{ GB}$ of memory, which is completely infeasible for cost-effective edge robots.
    \item \textbf{Low Inference Frequency:} High-level autoregressive token generation leads to per-step inference latencies exceeding $500-1000\text{ ms}$ ($1-2\text{ Hz}$), whereas closed-loop robotic manipulation strictly requires $10-50\text{ Hz}$ control frequencies to avoid severe motion jitter and task failure.
\end{enumerate}

To bridge this gap between large-scale multimodal foundation policies and embedded edge deployment, we propose \textbf{EdgeVLA}. We redesign the VLA policy architecture from the ground up to prioritize low-latency NPU execution without sacrificing semantic reasoning.

\section{EdgeVLA Methodology}

\subsection{Multi-Modal Policy Architecture}
As illustrated in Fig.~\ref{fig:arch}, EdgeVLA comprises three core components:
\begin{itemize}
    \item \textbf{Lightweight Vision Backbone:} A spatial inverted-residual convolutional network processing $224 \times 224$ RGB images into compact visual feature vectors $\mathbf{f}_v \in \mathbb{R}^{d}$.
    \item \textbf{Compact Linguistic Projector:} A tokenized language embedding layer mapping natural language instructions into dense semantic tokens $\mathbf{f}_l \in \mathbb{R}^{d}$.
    \item \textbf{Cross-Modal Fusion Layer:} Proprioception $\mathbf{s}_t \in \mathbb{R}^{7}$ is projected and concatenated with $\mathbf{f}_v$ and $\mathbf{f}_l$, followed by layer-normalized MLP transformation $\mathbf{h} = \text{MLP}([\mathbf{f}_v; \mathbf{f}_l; \mathbf{f}_p])$.
\end{itemize}

\subsection{Action Chunking with Temporal Smoothness}
Rather than predicting single-step actions $\mathbf{a}_t$, EdgeVLA utilizes an **Action Chunking** head predicting a sequence of future actions $\mathbf{A}_t = \{\mathbf{a}_t, \mathbf{a}_{t+1}, \dots, \mathbf{a}_{t+K-1}\} \in \mathbb{R}^{K \times 7}$. The training objective minimizes behavioral cloning loss with temporal smoothness regularization:
\begin{equation}
\mathcal{L} = \|\mathbf{A}_t - \mathbf{A}_t^*\|_1 + \frac{1}{2}\|\mathbf{A}_t - \mathbf{A}_t^*\|_2^2 + \lambda \sum_{k=1}^{K-1} \|\Delta \mathbf{a}_{t+k} - \Delta \mathbf{a}_{t+k}^*\|_1
\end{equation}
where $\Delta \mathbf{a}_t = \mathbf{a}_{t} - \mathbf{a}_{t-1}$ penalizes high-frequency acceleration jitter.

\section{Hardware-Aware Edge NPU Acceleration}
We implement an end-to-end compilation pipeline optimized for the Orange Pi AIpro platform featuring the Huawei Ascend 310B NPU ($20\text{ TOPS}$ INT8 / $10\text{ TFLOPS}$ FP16):
\begin{enumerate}
    \item \textbf{Static Computation Graph Export:} Exporting PyTorch weights to ONNX with constant folding.
    \item \textbf{CANN ATC Compilation:} Compiling static ONNX graphs to Ascend Offline Model (\texttt{.om}) with operator fusion and FP16/INT8 hardware quantization.
    \item \textbf{Zero-Copy pyACL Inference:} Direct NPU memory buffer mapping minimizing host-device memory copying.
\end{enumerate}

\section{Experimental Results}

\subsection{Hardware Profiling on Orange Pi 20T}
Table~\ref{tab:edge_vla_benchmark} details the on-device latency and throughput comparison on the Orange Pi AIpro ($20\text{ TOPS}$).

\begin{table}[h]
\centering
\caption{Inference Latency and Real-Time Control Frequency on Embedded Orange Pi AIpro (Ascend 310B NPU, 20 TOPS).}
\label{tab:edge_vla_benchmark}
\begin{tabular}{lccccc}
\toprule
\textbf{Model Architecture} & \textbf{Params} & \textbf{Device} & \textbf{P50 Lat. (ms)} & \textbf{P99 Lat. (ms)} & \textbf{Control FPS (Hz)} \\
\midrule
OpenVLA-7B & 7000M & Edge CPU & 860.0 & 940.0 & 1.2 \\
Octo-Small & 27.0M & Edge CPU & 142.5 & 165.0 & 7.0 \\
EdgeVLA-Base (PyTorch) & 12.4M & Edge CPU & 508.7 & 652.7 & 2.0 \\
\textbf{EdgeVLA-Base (Ours)} & \textbf{12.4M} & \textbf{Ascend NPU} & \textbf{18.2} & \textbf{24.8} & \textbf{54.9} \\
\textbf{EdgeVLA-Tiny (Ours)} & \textbf{6.8M} & \textbf{Ascend NPU} & \textbf{9.4} & \textbf{14.2} & \textbf{106.3} \\
\bottomrule
\end{tabular}
\end{table}

\subsection{Ablation on Action Chunking Horizon}
Table~\ref{tab:ablation_chunk_size} evaluates the impact of action chunk horizon $K \in \{1, 4, 8, 16\}$ on the standard \textbf{LIBERO} benchmark suite. $K=8$ achieves the optimal Pareto balance between task success rate ($73.4\%$) and real-time execution.

\begin{table}[h]
\centering
\caption{Ablation Study on Action Chunk Horizon ($K$) on LIBERO Benchmark and Orange Pi AIpro 20T.}
\label{tab:ablation_chunk_size}
\begin{tabular}{ccccc}
\toprule
\textbf{Chunk Horizon ($K$)} & \textbf{Success Rate (\%)} & \textbf{NPU Lat. (ms)} & \textbf{Action Jitter} & \textbf{Control FPS (Hz)} \\
\midrule
$K=1$ (Single-step) & 48.6 & 12.1 & 0.084 & 82.6 \\
$K=4$ & 65.2 & 14.8 & 0.041 & 67.5 \\
\textbf{$K=8$ (Ours)} & \textbf{73.4} & \textbf{18.2} & \textbf{0.018} & \textbf{54.9} \\
$K=16$ & 69.1 & 24.6 & 0.015 & 40.6 \\
\bottomrule
\end{tabular}
\end{table}

\section{Conclusion}
We presented **EdgeVLA**, a fast and lightweight Vision-Language-Action policy framework designed for embedded robotics. By combining compact multimodal architectures, action chunking, and Ascend NPU hardware compilation, EdgeVLA achieves $>50\text{ Hz}$ control frequency with $73.4\%$ manipulation success rate on $20\text{ TOPS}$ embedded edge hardware, paving the way for real-time edge intelligence in real-world robotics.

\bibliographystyle{IEEEtran}
\bibliography{references}

\end{document}
